\chapter{Planificación y cronograma}\label{anexo-f.-planificacion-y-esfuerzo}

Este anexo resume la planificación temporal del trabajo. Dado que no se mantuvo un registro horario exacto por actividad, en lugar de un desglose de horas se presenta un cronograma que refleja cuándo se desarrolló cada fase, que es la información de la que se dispone con fiabilidad. El desarrollo se concentró entre abril y junio de 2026.

Conviene aclarar una particularidad del proyecto: la obtención de los escenarios mediante la VPN universitaria se extendió durante varias semanas, pero fue un proceso en gran parte \emph{desatendido} (descargas que avanzaban en segundo plano y que solo requerían supervisión puntual para reanudar cortes de conexión y verificar la integridad de los ficheros frente al ETL). Por ello aparece en el cronograma como una franja larga de calendario solapada con el resto de tareas, y no como una dedicación equivalente en esfuerzo: el trabajo efectivo se concentró en las demás fases, que sí exigieron atención continua. Cuando las descargas se volvieron el cuello de botella, se sustituyeron por el transporte físico en disco duro (véase la sección~\ref{conclusiones-y-trabajo-futuro}).

\begin{table}[ht]
\centering
\setlength{\tabcolsep}{2pt}
\renewcommand{\arraystretch}{1.35}
\newcommand{\barra}{\textcolor{blue!55}{\rule{0.34cm}{7pt}}}
\small
\begin{tabular}{>{\raggedright\arraybackslash}p{5.1cm}*{13}{>{\centering\arraybackslash}m{0.34cm}}}
\toprule
 & \multicolumn{4}{c}{\textbf{Abril}} & \multicolumn{5}{c}{\textbf{Mayo}} & \multicolumn{4}{c}{\textbf{Junio}} \\
\cmidrule(lr){2-5}\cmidrule(lr){6-10}\cmidrule(lr){11-14}
\textbf{Fase} & 1 & 2 & 3 & 4 & 1 & 2 & 3 & 4 & 5 & 1 & 2 & 3 & 4 \\
\midrule
Formación e investigación        & \barra & \barra &       &       &       &       &       &       &       &       &       &       &       \\
Obtención de escenarios (VPN, desatendida) & \barra & \barra & \barra & \barra & \barra & \barra & \barra &       &       &       &       &       &       \\
Proceso ETL y \textit{benchmark} de almacenamiento &       & \barra & \barra & \barra &       &       &       &       &       &       &       &       &       \\
Motor de consultas analíticas    &       &       &       & \barra & \barra & \barra & \barra &       &       &       &       &       &       \\
Aplicación web                   &       &       &       &       &       & \barra & \barra & \barra & \barra &       &       &       &       \\
Despliegue y operación del servidor &     &       &       &       &       &       &       & \barra & \barra &       &       &       &       \\
Validación y pruebas             &       &       &       &       &       &       &       &       & \barra & \barra & \barra &       &       \\
Ampliación a los 56 escenarios   &       &       &       &       &       &       &       &       &       & \barra & \barra &       &       \\
Evaluación experimental multi-máquina &  &       &       &       &       &       &       &       &       &       & \barra & \barra &       \\
Integración con sistemas GIS     &       &       &       &       &       &       &       &       &       &       &       & \barra &       \\
Redacción de la memoria          &       &       &       &       &       &       &       & \barra & \barra & \barra & \barra & \barra & \barra \\
\bottomrule
\end{tabular}
\caption{Cronograma del proyecto por semanas (abril--junio de 2026).}
\label{tab:cronograma}
\end{table}

Las fases no fueron estrictamente secuenciales. El ETL y el motor de consultas se revisaron varias veces a medida que la validación y el uso real de la aplicación revelaban problemas de rendimiento o de memoria, y la redacción de la memoria se solapó con las últimas etapas de desarrollo. Las dos últimas fases del cronograma ---la evaluación del rendimiento en distintas máquinas (sección~\ref{rendimiento-multimaquina}) y la interoperabilidad con sistemas GIS (capítulo~\ref{cap:gis})--- ampliaron el alcance inicial una vez consolidado el núcleo del sistema.
